\documentclass{article}
\usepackage{amsmath, amssymb, amsthm}
\usepackage{hyperref}
\usepackage{geometry}
\geometry{margin=1in}

\title{Erd\H{o}s Problem \#323}
\date{Last updated: 2025-08-31}
\author{Source: erdosproblems.com}

\begin{document}
\maketitle

\noindent\textbf{Status:} OPEN \quad \textbf{No prize}

\noindent\textbf{Tags:} number theory, powers

\medskip
\noindent\textbf{Problem Statement:}

Let $1\leq m\leq k$ and $f_{k,m}(x)$ denote the number of integers $\leq x$ which are the sum of $m$ many nonnegative $k$th powers. Is it true that\[f_{k,k}(x) \gg_\epsilon x^{1-\epsilon}\]for all $\epsilon>0$? Is it true that if $m<k$ then\[f_{k,m}(x) \gg x^{m/k}\]for sufficiently large $x$?

\bigskip
\noindent\textbf{Remarks:}

This would have significant applications to Waring's problem. Erd\H{o}s and Graham describe this as 'unattackable by the methods at our disposal'. The case $k=2$ was resolved by Landau, who showed\[f_{2,2}(x) \sim \frac{cx}{\sqrt{\log x}}\]for some constant $c>0$.

For $k>2$ it is not known if $f_{k,k}(x)=o(x)$.

\bigskip
\noindent\small{Source: \url{https://www.erdosproblems.com/323}}
\end{document}
